\chapter*{Introduction}
\addcontentsline{toc}{chapter}{Introduction}

La retinopatía diabética (RD) es una de las principales causas de ceguera irreversible en el mundo. Su progresión suele ser silenciosa y asintomática hasta estadios avanzados, por lo que el diagnóstico temprano resulta clave para prevenir complicaciones severas. Sin embargo, en contextos con recursos limitados, como centros de atención primaria, áreas rurales o dispositivos móviles de tamizaje, el acceso al diagnóstico oftalmológico especializado suele estar restringido.

En este escenario, muchas estrategias de detección temprana dependen de la captura de imágenes retinianas por parte de técnicos no especialistas. Estas imágenes son luego evaluadas por profesionales o sistemas automatizados. En ese flujo de trabajo, contar con mecanismos que garanticen la calidad de las imágenes capturadas es esencial: imágenes desenfocadas, mal iluminadas o afectadas por artefactos no solo dificultan la interpretación clínica, sino que obligan a repetir el estudio, con consecuencias negativas para los pacientes y el sistema de salud.

Frente a este desafío, la presente tesis se propone desarrollar una herramienta de control de calidad automatizado para imágenes de fondo de ojo, capaz de clasificarlas como de buena o mala calidad en el momento de su captura. Esta herramienta busca integrarse al proceso de tamizaje, permitiendo alertar al operador si una imagen no es apta para diagnóstico, antes de que el paciente abandone la consulta.

Desde el punto de vista técnico, gran parte de la bibliografía en este campo se basa en el entrenamiento de modelos convolucionales supervisados, específicamente diseñados y ajustados para esta tarea. Este enfoque, si bien efectivo, implica un esfuerzo considerable en términos de recolección de datos, curación, etiquetado, ajuste de hiperparámetros y entrenamiento.

En contraste con este paradigma, en los últimos años han surgido modelos multimodales de propósito general, como GPT-4o, capaces de procesar imágenes y texto en simultáneo. Estos modelos permiten realizar tareas complejas a partir de instrucciones en lenguaje natural, sin necesidad de entrenamiento adicional específico. Esta tesis propone comparar ambos enfoques: el modelo entrenado (una ResNet-18 ajustada mediante transferencia de aprendizaje) frente a un modelo no entrenado (zero-shot y few-shot), evaluando sus fortalezas, limitaciones y posibles aplicaciones en contextos reales.

Esta comparación busca no solo medir el desempeño de ambos modelos, sino también reflexionar sobre el cambio de paradigma que representan las soluciones de bajo esfuerzo basadas en modelos multimodales, y su potencial para democratizar el acceso a herramientas de diagnóstico asistido por inteligencia artificial.

En este marco, los objetivos de la tesis son: (1) desarrollar un modelo entrenado específicamente para la clasificación de calidad de imágenes retinianas mediante una red ResNet-18 fine-tuned, (2) evaluar el desempeño de GPT-4o sin entrenamiento adicional, en modalidades zero-shot y few-shot, y (3) comparar ambos enfoques en términos de precisión, robustez, interpretabilidad y esfuerzo requerido para su implementación. El trabajo se organiza en los siguientes capítulos: en primer lugar, se presenta el marco teórico y los antecedentes relevantes en el área; luego, se detallan la metodología empleada y las características del conjunto de datos utilizado; a continuación, se describen los modelos implementados y el diseño experimental; posteriormente, se analizan y discuten los resultados obtenidos; finalmente, se presentan las conclusiones y se discuten las implicancias del trabajo en el contexto clínico y tecnológico actual.
